% !TeX program = pdflatex
% !TeX TXS-program:compile = txs:///pdflatex/[-shell-escape]
% !TeX TXS-program:pdflatex = pdflatex -synctex=1 -interaction=nonstopmode --shell-escape %.tex

\documentclass[../../../main_compilation.tex]{subfiles}

\begin{document}

% ================================================================================
% Encabezado del Reto CTF
% ================================================================================
% Nota: Si tienes un logotipo o imagen específica para este CTF, puedes definirlo con:
% \setCustomHeader{figure/logo_ctf.png}{NOMBRE_RETO}{\today}
% O usar el icono por defecto de CTF a través de \targetSummary:

\section[CTF - NOMBRE\_RETO (Evento)]{NOMBRE\_RETO — Evento CTF}
\label{sec:ctf-nombre-reto}

\targetSummary%
  {NOMBRE\_RETO}% Nombre del reto (Aparece en la cabecera derecha)
  {challenge.ctf.io:1337}% Host o archivo adjunto
  {CTF}% Plataforma / Ícono (CTF, HackTheBox, TryHackMe, o ruta de imagen 'figure/logo.png')
  {Medium}% Dificultad: Easy, Medium, Hard, Insane
  {Binary / Web / Crypto / Reversing}% Categoría
  {Puntos: 500}% Puntos o Tags
  {\today}% Fecha de resolución manual (Aparece en el pie de página izquierdo)
  {Solucionado}% Estado
  {flag\{...root\_flag...\}}% Flag

\vspace{0.4cm}

% ================================================================================
% 1. Enunciado del Reto
% ================================================================================
\subsection{Enunciado y Archivos Provistos}
Descripción del reto dada por los organizadores del CTF y lista de archivos entregados (binario, volcado de memoria, captura \texttt{.pcap}, script, etc.).

% ================================================================================
% 2. Análisis y Solución (Solver Script)
% ================================================================================
\subsection{Análisis \& Desarrollo del Exploit / Solver}

\begin{terminal}[kali@pentest:\~{}]
# Comandos de análisis estático / dinámico (GDB, Ghidra, Volatility, Wireshark, etc.)
checksec --file=vuln_binary
file vuln_binary
\end{terminal}

\begin{codebox}[Python]{solve.py}
from pwn import *

# Solver script automatizado
# context.log_level = 'debug'
# io = remote('challenge.ctf.io', 1337)
# io.sendline(...)
# print(io.recvall().decode())
\end{codebox}

\begin{flagbox}{CTF Flag}
flag{c0ngr4ts_y0u_s0lv3d_th1s_ch4ll3ng3}
\end{flagbox}

% ================================================================================
% 3. Conceptos y Técnicas Aprendidas
% ================================================================================
\subsection{Conocimientos y Técnicas Aprendidas}

\begin{learningbox}[Lecciones del Reto CTF]
	\begin{itemize}[leftmargin=*]
		\item \textbf{Técnica / Concepto Pwn/Crypto/Reversing:} Mecanismo matemático o vulnerabilidad explotada.
		\item \textbf{Herramienta / Biblioteca Clave:} Librerías (e.g. \texttt{pwntools}, \texttt{z3}, \texttt{angr}, Ghidra, GDB-gef).
	\end{itemize}
\end{learningbox}

\begin{sourcebox}[Fuentes y Referencias del Desafío]
	\begin{itemize}[leftmargin=*]
		\resourceItem{Web}{CTFtime / Writeups Oficiales}{\url{https://ctftime.org}}
		\resourceItem{Docs}{Documentación de Herramientas (pwntools / GEF)}{\url{https://docs.pwntools.com}}
	\end{itemize}
\end{sourcebox}

\end{document}
